
% Default to the notebook output style

    


% Inherit from the specified cell style.




    
\documentclass[11pt]{article}

    
    
    \usepackage[T1]{fontenc}
    % Nicer default font (+ math font) than Computer Modern for most use cases
    \usepackage{mathpazo}

    % Basic figure setup, for now with no caption control since it's done
    % automatically by Pandoc (which extracts ![](path) syntax from Markdown).
    \usepackage{graphicx}
    % We will generate all images so they have a width \maxwidth. This means
    % that they will get their normal width if they fit onto the page, but
    % are scaled down if they would overflow the margins.
    \makeatletter
    \def\maxwidth{\ifdim\Gin@nat@width>\linewidth\linewidth
    \else\Gin@nat@width\fi}
    \makeatother
    \let\Oldincludegraphics\includegraphics
    % Set max figure width to be 80% of text width, for now hardcoded.
    \renewcommand{\includegraphics}[1]{\Oldincludegraphics[width=.8\maxwidth]{#1}}
    % Ensure that by default, figures have no caption (until we provide a
    % proper Figure object with a Caption API and a way to capture that
    % in the conversion process - todo).
    \usepackage{caption}
    \DeclareCaptionLabelFormat{nolabel}{}
    \captionsetup{labelformat=nolabel}

    \usepackage{adjustbox} % Used to constrain images to a maximum size 
    \usepackage{xcolor} % Allow colors to be defined
    \usepackage{enumerate} % Needed for markdown enumerations to work
    \usepackage{geometry} % Used to adjust the document margins
    \usepackage{amsmath} % Equations
    \usepackage{amssymb} % Equations
    \usepackage{textcomp} % defines textquotesingle
    % Hack from http://tex.stackexchange.com/a/47451/13684:
    \AtBeginDocument{%
        \def\PYZsq{\textquotesingle}% Upright quotes in Pygmentized code
    }
    \usepackage{upquote} % Upright quotes for verbatim code
    \usepackage{eurosym} % defines \euro
    \usepackage[mathletters]{ucs} % Extended unicode (utf-8) support
    \usepackage[utf8x]{inputenc} % Allow utf-8 characters in the tex document
    \usepackage{fancyvrb} % verbatim replacement that allows latex
    \usepackage{grffile} % extends the file name processing of package graphics 
                         % to support a larger range 
    % The hyperref package gives us a pdf with properly built
    % internal navigation ('pdf bookmarks' for the table of contents,
    % internal cross-reference links, web links for URLs, etc.)
    \usepackage{hyperref}
    \usepackage{longtable} % longtable support required by pandoc >1.10
    \usepackage{booktabs}  % table support for pandoc > 1.12.2
    \usepackage[inline]{enumitem} % IRkernel/repr support (it uses the enumerate* environment)
    \usepackage[normalem]{ulem} % ulem is needed to support strikethroughs (\sout)
                                % normalem makes italics be italics, not underlines
    

    
    
    % Colors for the hyperref package
    \definecolor{urlcolor}{rgb}{0,.145,.698}
    \definecolor{linkcolor}{rgb}{.71,0.21,0.01}
    \definecolor{citecolor}{rgb}{.12,.54,.11}

    % ANSI colors
    \definecolor{ansi-black}{HTML}{3E424D}
    \definecolor{ansi-black-intense}{HTML}{282C36}
    \definecolor{ansi-red}{HTML}{E75C58}
    \definecolor{ansi-red-intense}{HTML}{B22B31}
    \definecolor{ansi-green}{HTML}{00A250}
    \definecolor{ansi-green-intense}{HTML}{007427}
    \definecolor{ansi-yellow}{HTML}{DDB62B}
    \definecolor{ansi-yellow-intense}{HTML}{B27D12}
    \definecolor{ansi-blue}{HTML}{208FFB}
    \definecolor{ansi-blue-intense}{HTML}{0065CA}
    \definecolor{ansi-magenta}{HTML}{D160C4}
    \definecolor{ansi-magenta-intense}{HTML}{A03196}
    \definecolor{ansi-cyan}{HTML}{60C6C8}
    \definecolor{ansi-cyan-intense}{HTML}{258F8F}
    \definecolor{ansi-white}{HTML}{C5C1B4}
    \definecolor{ansi-white-intense}{HTML}{A1A6B2}

    % commands and environments needed by pandoc snippets
    % extracted from the output of `pandoc -s`
    \providecommand{\tightlist}{%
      \setlength{\itemsep}{0pt}\setlength{\parskip}{0pt}}
    \DefineVerbatimEnvironment{Highlighting}{Verbatim}{commandchars=\\\{\}}
    % Add ',fontsize=\small' for more characters per line
    \newenvironment{Shaded}{}{}
    \newcommand{\KeywordTok}[1]{\textcolor[rgb]{0.00,0.44,0.13}{\textbf{{#1}}}}
    \newcommand{\DataTypeTok}[1]{\textcolor[rgb]{0.56,0.13,0.00}{{#1}}}
    \newcommand{\DecValTok}[1]{\textcolor[rgb]{0.25,0.63,0.44}{{#1}}}
    \newcommand{\BaseNTok}[1]{\textcolor[rgb]{0.25,0.63,0.44}{{#1}}}
    \newcommand{\FloatTok}[1]{\textcolor[rgb]{0.25,0.63,0.44}{{#1}}}
    \newcommand{\CharTok}[1]{\textcolor[rgb]{0.25,0.44,0.63}{{#1}}}
    \newcommand{\StringTok}[1]{\textcolor[rgb]{0.25,0.44,0.63}{{#1}}}
    \newcommand{\CommentTok}[1]{\textcolor[rgb]{0.38,0.63,0.69}{\textit{{#1}}}}
    \newcommand{\OtherTok}[1]{\textcolor[rgb]{0.00,0.44,0.13}{{#1}}}
    \newcommand{\AlertTok}[1]{\textcolor[rgb]{1.00,0.00,0.00}{\textbf{{#1}}}}
    \newcommand{\FunctionTok}[1]{\textcolor[rgb]{0.02,0.16,0.49}{{#1}}}
    \newcommand{\RegionMarkerTok}[1]{{#1}}
    \newcommand{\ErrorTok}[1]{\textcolor[rgb]{1.00,0.00,0.00}{\textbf{{#1}}}}
    \newcommand{\NormalTok}[1]{{#1}}
    
    % Additional commands for more recent versions of Pandoc
    \newcommand{\ConstantTok}[1]{\textcolor[rgb]{0.53,0.00,0.00}{{#1}}}
    \newcommand{\SpecialCharTok}[1]{\textcolor[rgb]{0.25,0.44,0.63}{{#1}}}
    \newcommand{\VerbatimStringTok}[1]{\textcolor[rgb]{0.25,0.44,0.63}{{#1}}}
    \newcommand{\SpecialStringTok}[1]{\textcolor[rgb]{0.73,0.40,0.53}{{#1}}}
    \newcommand{\ImportTok}[1]{{#1}}
    \newcommand{\DocumentationTok}[1]{\textcolor[rgb]{0.73,0.13,0.13}{\textit{{#1}}}}
    \newcommand{\AnnotationTok}[1]{\textcolor[rgb]{0.38,0.63,0.69}{\textbf{\textit{{#1}}}}}
    \newcommand{\CommentVarTok}[1]{\textcolor[rgb]{0.38,0.63,0.69}{\textbf{\textit{{#1}}}}}
    \newcommand{\VariableTok}[1]{\textcolor[rgb]{0.10,0.09,0.49}{{#1}}}
    \newcommand{\ControlFlowTok}[1]{\textcolor[rgb]{0.00,0.44,0.13}{\textbf{{#1}}}}
    \newcommand{\OperatorTok}[1]{\textcolor[rgb]{0.40,0.40,0.40}{{#1}}}
    \newcommand{\BuiltInTok}[1]{{#1}}
    \newcommand{\ExtensionTok}[1]{{#1}}
    \newcommand{\PreprocessorTok}[1]{\textcolor[rgb]{0.74,0.48,0.00}{{#1}}}
    \newcommand{\AttributeTok}[1]{\textcolor[rgb]{0.49,0.56,0.16}{{#1}}}
    \newcommand{\InformationTok}[1]{\textcolor[rgb]{0.38,0.63,0.69}{\textbf{\textit{{#1}}}}}
    \newcommand{\WarningTok}[1]{\textcolor[rgb]{0.38,0.63,0.69}{\textbf{\textit{{#1}}}}}
    
    
    % Define a nice break command that doesn't care if a line doesn't already
    % exist.
    \def\br{\hspace*{\fill} \\* }
    % Math Jax compatability definitions
    \def\gt{>}
    \def\lt{<}
    % Document parameters
    \title{Numerical study of the main theorem (sd=0.020, beneficial=1e-14)}
    
    
    

    % Pygments definitions
    
\makeatletter
\def\PY@reset{\let\PY@it=\relax \let\PY@bf=\relax%
    \let\PY@ul=\relax \let\PY@tc=\relax%
    \let\PY@bc=\relax \let\PY@ff=\relax}
\def\PY@tok#1{\csname PY@tok@#1\endcsname}
\def\PY@toks#1+{\ifx\relax#1\empty\else%
    \PY@tok{#1}\expandafter\PY@toks\fi}
\def\PY@do#1{\PY@bc{\PY@tc{\PY@ul{%
    \PY@it{\PY@bf{\PY@ff{#1}}}}}}}
\def\PY#1#2{\PY@reset\PY@toks#1+\relax+\PY@do{#2}}

\expandafter\def\csname PY@tok@w\endcsname{\def\PY@tc##1{\textcolor[rgb]{0.73,0.73,0.73}{##1}}}
\expandafter\def\csname PY@tok@c\endcsname{\let\PY@it=\textit\def\PY@tc##1{\textcolor[rgb]{0.25,0.50,0.50}{##1}}}
\expandafter\def\csname PY@tok@cp\endcsname{\def\PY@tc##1{\textcolor[rgb]{0.74,0.48,0.00}{##1}}}
\expandafter\def\csname PY@tok@k\endcsname{\let\PY@bf=\textbf\def\PY@tc##1{\textcolor[rgb]{0.00,0.50,0.00}{##1}}}
\expandafter\def\csname PY@tok@kp\endcsname{\def\PY@tc##1{\textcolor[rgb]{0.00,0.50,0.00}{##1}}}
\expandafter\def\csname PY@tok@kt\endcsname{\def\PY@tc##1{\textcolor[rgb]{0.69,0.00,0.25}{##1}}}
\expandafter\def\csname PY@tok@o\endcsname{\def\PY@tc##1{\textcolor[rgb]{0.40,0.40,0.40}{##1}}}
\expandafter\def\csname PY@tok@ow\endcsname{\let\PY@bf=\textbf\def\PY@tc##1{\textcolor[rgb]{0.67,0.13,1.00}{##1}}}
\expandafter\def\csname PY@tok@nb\endcsname{\def\PY@tc##1{\textcolor[rgb]{0.00,0.50,0.00}{##1}}}
\expandafter\def\csname PY@tok@nf\endcsname{\def\PY@tc##1{\textcolor[rgb]{0.00,0.00,1.00}{##1}}}
\expandafter\def\csname PY@tok@nc\endcsname{\let\PY@bf=\textbf\def\PY@tc##1{\textcolor[rgb]{0.00,0.00,1.00}{##1}}}
\expandafter\def\csname PY@tok@nn\endcsname{\let\PY@bf=\textbf\def\PY@tc##1{\textcolor[rgb]{0.00,0.00,1.00}{##1}}}
\expandafter\def\csname PY@tok@ne\endcsname{\let\PY@bf=\textbf\def\PY@tc##1{\textcolor[rgb]{0.82,0.25,0.23}{##1}}}
\expandafter\def\csname PY@tok@nv\endcsname{\def\PY@tc##1{\textcolor[rgb]{0.10,0.09,0.49}{##1}}}
\expandafter\def\csname PY@tok@no\endcsname{\def\PY@tc##1{\textcolor[rgb]{0.53,0.00,0.00}{##1}}}
\expandafter\def\csname PY@tok@nl\endcsname{\def\PY@tc##1{\textcolor[rgb]{0.63,0.63,0.00}{##1}}}
\expandafter\def\csname PY@tok@ni\endcsname{\let\PY@bf=\textbf\def\PY@tc##1{\textcolor[rgb]{0.60,0.60,0.60}{##1}}}
\expandafter\def\csname PY@tok@na\endcsname{\def\PY@tc##1{\textcolor[rgb]{0.49,0.56,0.16}{##1}}}
\expandafter\def\csname PY@tok@nt\endcsname{\let\PY@bf=\textbf\def\PY@tc##1{\textcolor[rgb]{0.00,0.50,0.00}{##1}}}
\expandafter\def\csname PY@tok@nd\endcsname{\def\PY@tc##1{\textcolor[rgb]{0.67,0.13,1.00}{##1}}}
\expandafter\def\csname PY@tok@s\endcsname{\def\PY@tc##1{\textcolor[rgb]{0.73,0.13,0.13}{##1}}}
\expandafter\def\csname PY@tok@sd\endcsname{\let\PY@it=\textit\def\PY@tc##1{\textcolor[rgb]{0.73,0.13,0.13}{##1}}}
\expandafter\def\csname PY@tok@si\endcsname{\let\PY@bf=\textbf\def\PY@tc##1{\textcolor[rgb]{0.73,0.40,0.53}{##1}}}
\expandafter\def\csname PY@tok@se\endcsname{\let\PY@bf=\textbf\def\PY@tc##1{\textcolor[rgb]{0.73,0.40,0.13}{##1}}}
\expandafter\def\csname PY@tok@sr\endcsname{\def\PY@tc##1{\textcolor[rgb]{0.73,0.40,0.53}{##1}}}
\expandafter\def\csname PY@tok@ss\endcsname{\def\PY@tc##1{\textcolor[rgb]{0.10,0.09,0.49}{##1}}}
\expandafter\def\csname PY@tok@sx\endcsname{\def\PY@tc##1{\textcolor[rgb]{0.00,0.50,0.00}{##1}}}
\expandafter\def\csname PY@tok@m\endcsname{\def\PY@tc##1{\textcolor[rgb]{0.40,0.40,0.40}{##1}}}
\expandafter\def\csname PY@tok@gh\endcsname{\let\PY@bf=\textbf\def\PY@tc##1{\textcolor[rgb]{0.00,0.00,0.50}{##1}}}
\expandafter\def\csname PY@tok@gu\endcsname{\let\PY@bf=\textbf\def\PY@tc##1{\textcolor[rgb]{0.50,0.00,0.50}{##1}}}
\expandafter\def\csname PY@tok@gd\endcsname{\def\PY@tc##1{\textcolor[rgb]{0.63,0.00,0.00}{##1}}}
\expandafter\def\csname PY@tok@gi\endcsname{\def\PY@tc##1{\textcolor[rgb]{0.00,0.63,0.00}{##1}}}
\expandafter\def\csname PY@tok@gr\endcsname{\def\PY@tc##1{\textcolor[rgb]{1.00,0.00,0.00}{##1}}}
\expandafter\def\csname PY@tok@ge\endcsname{\let\PY@it=\textit}
\expandafter\def\csname PY@tok@gs\endcsname{\let\PY@bf=\textbf}
\expandafter\def\csname PY@tok@gp\endcsname{\let\PY@bf=\textbf\def\PY@tc##1{\textcolor[rgb]{0.00,0.00,0.50}{##1}}}
\expandafter\def\csname PY@tok@go\endcsname{\def\PY@tc##1{\textcolor[rgb]{0.53,0.53,0.53}{##1}}}
\expandafter\def\csname PY@tok@gt\endcsname{\def\PY@tc##1{\textcolor[rgb]{0.00,0.27,0.87}{##1}}}
\expandafter\def\csname PY@tok@err\endcsname{\def\PY@bc##1{\setlength{\fboxsep}{0pt}\fcolorbox[rgb]{1.00,0.00,0.00}{1,1,1}{\strut ##1}}}
\expandafter\def\csname PY@tok@kc\endcsname{\let\PY@bf=\textbf\def\PY@tc##1{\textcolor[rgb]{0.00,0.50,0.00}{##1}}}
\expandafter\def\csname PY@tok@kd\endcsname{\let\PY@bf=\textbf\def\PY@tc##1{\textcolor[rgb]{0.00,0.50,0.00}{##1}}}
\expandafter\def\csname PY@tok@kn\endcsname{\let\PY@bf=\textbf\def\PY@tc##1{\textcolor[rgb]{0.00,0.50,0.00}{##1}}}
\expandafter\def\csname PY@tok@kr\endcsname{\let\PY@bf=\textbf\def\PY@tc##1{\textcolor[rgb]{0.00,0.50,0.00}{##1}}}
\expandafter\def\csname PY@tok@bp\endcsname{\def\PY@tc##1{\textcolor[rgb]{0.00,0.50,0.00}{##1}}}
\expandafter\def\csname PY@tok@fm\endcsname{\def\PY@tc##1{\textcolor[rgb]{0.00,0.00,1.00}{##1}}}
\expandafter\def\csname PY@tok@vc\endcsname{\def\PY@tc##1{\textcolor[rgb]{0.10,0.09,0.49}{##1}}}
\expandafter\def\csname PY@tok@vg\endcsname{\def\PY@tc##1{\textcolor[rgb]{0.10,0.09,0.49}{##1}}}
\expandafter\def\csname PY@tok@vi\endcsname{\def\PY@tc##1{\textcolor[rgb]{0.10,0.09,0.49}{##1}}}
\expandafter\def\csname PY@tok@vm\endcsname{\def\PY@tc##1{\textcolor[rgb]{0.10,0.09,0.49}{##1}}}
\expandafter\def\csname PY@tok@sa\endcsname{\def\PY@tc##1{\textcolor[rgb]{0.73,0.13,0.13}{##1}}}
\expandafter\def\csname PY@tok@sb\endcsname{\def\PY@tc##1{\textcolor[rgb]{0.73,0.13,0.13}{##1}}}
\expandafter\def\csname PY@tok@sc\endcsname{\def\PY@tc##1{\textcolor[rgb]{0.73,0.13,0.13}{##1}}}
\expandafter\def\csname PY@tok@dl\endcsname{\def\PY@tc##1{\textcolor[rgb]{0.73,0.13,0.13}{##1}}}
\expandafter\def\csname PY@tok@s2\endcsname{\def\PY@tc##1{\textcolor[rgb]{0.73,0.13,0.13}{##1}}}
\expandafter\def\csname PY@tok@sh\endcsname{\def\PY@tc##1{\textcolor[rgb]{0.73,0.13,0.13}{##1}}}
\expandafter\def\csname PY@tok@s1\endcsname{\def\PY@tc##1{\textcolor[rgb]{0.73,0.13,0.13}{##1}}}
\expandafter\def\csname PY@tok@mb\endcsname{\def\PY@tc##1{\textcolor[rgb]{0.40,0.40,0.40}{##1}}}
\expandafter\def\csname PY@tok@mf\endcsname{\def\PY@tc##1{\textcolor[rgb]{0.40,0.40,0.40}{##1}}}
\expandafter\def\csname PY@tok@mh\endcsname{\def\PY@tc##1{\textcolor[rgb]{0.40,0.40,0.40}{##1}}}
\expandafter\def\csname PY@tok@mi\endcsname{\def\PY@tc##1{\textcolor[rgb]{0.40,0.40,0.40}{##1}}}
\expandafter\def\csname PY@tok@il\endcsname{\def\PY@tc##1{\textcolor[rgb]{0.40,0.40,0.40}{##1}}}
\expandafter\def\csname PY@tok@mo\endcsname{\def\PY@tc##1{\textcolor[rgb]{0.40,0.40,0.40}{##1}}}
\expandafter\def\csname PY@tok@ch\endcsname{\let\PY@it=\textit\def\PY@tc##1{\textcolor[rgb]{0.25,0.50,0.50}{##1}}}
\expandafter\def\csname PY@tok@cm\endcsname{\let\PY@it=\textit\def\PY@tc##1{\textcolor[rgb]{0.25,0.50,0.50}{##1}}}
\expandafter\def\csname PY@tok@cpf\endcsname{\let\PY@it=\textit\def\PY@tc##1{\textcolor[rgb]{0.25,0.50,0.50}{##1}}}
\expandafter\def\csname PY@tok@c1\endcsname{\let\PY@it=\textit\def\PY@tc##1{\textcolor[rgb]{0.25,0.50,0.50}{##1}}}
\expandafter\def\csname PY@tok@cs\endcsname{\let\PY@it=\textit\def\PY@tc##1{\textcolor[rgb]{0.25,0.50,0.50}{##1}}}

\def\PYZbs{\char`\\}
\def\PYZus{\char`\_}
\def\PYZob{\char`\{}
\def\PYZcb{\char`\}}
\def\PYZca{\char`\^}
\def\PYZam{\char`\&}
\def\PYZlt{\char`\<}
\def\PYZgt{\char`\>}
\def\PYZsh{\char`\#}
\def\PYZpc{\char`\%}
\def\PYZdl{\char`\$}
\def\PYZhy{\char`\-}
\def\PYZsq{\char`\'}
\def\PYZdq{\char`\"}
\def\PYZti{\char`\~}
% for compatibility with earlier versions
\def\PYZat{@}
\def\PYZlb{[}
\def\PYZrb{]}
\makeatother


    % Exact colors from NB
    \definecolor{incolor}{rgb}{0.0, 0.0, 0.5}
    \definecolor{outcolor}{rgb}{0.545, 0.0, 0.0}



    
    % Prevent overflowing lines due to hard-to-break entities
    \sloppy 
    % Setup hyperref package
    \hypersetup{
      breaklinks=true,  % so long urls are correctly broken across lines
      colorlinks=true,
      urlcolor=urlcolor,
      linkcolor=linkcolor,
      citecolor=citecolor,
      }
    % Slightly bigger margins than the latex defaults
    
    \geometry{verbose,tmargin=1in,bmargin=1in,lmargin=1in,rmargin=1in}
    
    

    \begin{document}
    
    
    \maketitle
    
    

    
    \section{Numerical Study of the Main Theorem (sd=0.020,
beneficial=1e-14)}\label{numerical-study-of-the-main-theorem-sd0.020-beneficial1e-14}

\textbf{Parameters.} There are nine versions of this notebook. The only
cell in which they differ is the following, where each of two parameters
is assigned one of three possible values (there are \(3 \times 3 = 9\)
joint assignments). The value of "sd" in the title should match that of
the first parameter, the standard deviation in fitness of the initial
population, and the value of "beneficial" should match that of the
second parameter, the fraction of mutations that increase fitness.

    \begin{Verbatim}[commandchars=\\\{\}]
{\color{incolor}In [{\color{incolor}1}]:} \PY{n}{standard\PYZus{}deviation\PYZus{}in\PYZus{}fitness\PYZus{}of\PYZus{}initial\PYZus{}population} \PY{o}{=} \PY{l+s+s1}{\PYZsq{}}\PY{l+s+s1}{0.020}\PY{l+s+s1}{\PYZsq{}} \PY{c+c1}{\PYZsh{} Levels: 0.006, 0.013, 0.020}
        \PY{n}{fraction\PYZus{}of\PYZus{}mutations\PYZus{}that\PYZus{}increase\PYZus{}fitness} \PY{o}{=} \PY{l+s+s1}{\PYZsq{}}\PY{l+s+s1}{1e\PYZhy{}14}\PY{l+s+s1}{\PYZsq{}}         \PY{c+c1}{\PYZsh{} Levels: 1e\PYZhy{}06, 1e\PYZhy{}10, 1e\PYZhy{}14}
        
        \PY{o}{!}\PY{n+nv}{TZ}\PY{o}{=}America/Los\PYZus{}Angeles date
\end{Verbatim}


    \begin{Verbatim}[commandchars=\\\{\}]
Fri Nov  2 04:11:56 PDT 2018

    \end{Verbatim}

    \subsection{Introduction}\label{introduction}

We investigate numerically
\href{https://link.springer.com/article/10.1007/s00285-017-1190-x\#Equ10}{Equation
3.4} in the main theorem of Basener and Sanford, which expresses the
instantaneous rate of change in mean fitness of the population as the
sum of two terms:

\begin{align*}
   \frac{\text{d}\bar{m}}{\text{d}t} 
      &= \text{Var}(m) 
         + \underbrace{\frac{1}{\sum_i P_i} \sum_i (B^\text{in}_i - B^\text{out}_i)(m_i - \bar{m})}_\text{correction for mutational effects}.
\end{align*}

The first term on the right-hand side is the variance in fitness of a
population with frequency distribution \(P\) over the fitness classes
\(m_i.\) In the absence of mutation, the instantaneous rate of change in
mean fitness \(\bar{m}\) of the population is equal to the variance
\(\text{Var}(m)\) of the population in fitness. The second term is a
correction for mutational effects, taking into account the differences
in rates \(B^\text{in}_i\) of birth \textbf{of} organisms in fitness
classes and the rates \(B^\text{out}_i\) of birth \textbf{to} organisms
in the classes. If there were no mutation, or if mutation never had an
effect on the fitness class of the offspring, then every birth to an
organism with fitness \(m_i\) would be a birth of an organism with
fitness \(m_i,\) and it would hold that
\(B_i^\text{in} - B_i^\text{out} = 0\) for all \(i.\)

The variance is of course non-negative, and is zero only when the
population is comprised entirely of organisms in a single fitness class.
The correction term can be positive early in the evolutionary process,
at least in contrived circumstances. However, as the process converges
to the equilibrium distribution of the population over fitness classes,
the rate of change in mean fitness goes to zero, and the correction
comes to equal the negative variance.

\subsubsection{Overview of the
experiment}\label{overview-of-the-experiment}

{[}{More here}{]}

When mutant offspring range over all fitness classes, as predicated by
Basener and Sanford in Sections 5.3 and 5.4 of their article,
\textbf{the equilibrium distribution \emph{does not} depend on the
initial distribution of the population.} The purpose in experimenting
with different settings of initial variance in fitness is to observe the
effect on evolutionary dynamics. There is no effect on the "ultimate
fate" of the population. The equilibrium distribution does depend on the
fraction of mutations that are beneficial. Unsurprisingly, an effect of
reduction in the fraction is reduction of the mean fitness at
equilibrium. Surprisingly, we can set the fraction as low as accurate
computation with conventional floating-point numbers (binary64)
permits~--- much smaller than Basener and Sanford indicate is
biologically realistic~--- and the mean fitness at equilibrium remains
positive.

\textbf{Note on viewing results.} We recommend opening different
versions of this notebook in different windows, and placing the windows
side-by-side. It may be helpful, when comparing graphical outputs, to
reduce the scale of display in the windows.

    \begin{Verbatim}[commandchars=\\\{\}]
{\color{incolor}In [{\color{incolor}2}]:} \PY{o}{\PYZpc{}}\PY{k}{matplotlib} notebook
        \PY{o}{\PYZpc{}}\PY{k}{run} ../../Code/bs.py
        \PY{o}{\PYZpc{}}\PY{k}{run} \PYZhy{}i ../../Code/multiprecision\PYZus{}gamma.py
        \PY{o}{\PYZpc{}}\PY{k}{run} \PYZhy{}i ../../Code/gamma.py
\end{Verbatim}


    \subsection{1~~Set up the numerical
experiment}\label{set-up-the-numerical-experiment}

{[}{Summarize steps}{]}

    \subsubsection{1.1~~Define fitness classes and related
parameters}\label{define-fitness-classes-and-related-parameters}

The fitnesses \(m_i\) are spaced uniformly over the interval {[}-0.1,
0.1{]}, with endpoints included. The spacing of 0.0005 comes from
Basener and Sanford (Section 5.4). Although they indicate that the
maximum possible fitness is 0.15, they effectively limit it to 0.1 in
the numerical experiments that they report. We explicitly impose the
limit, and thus avoid an objection to our methodology that Basener and
Sanford otherwise might have raised.

    \begin{Verbatim}[commandchars=\\\{\}]
{\color{incolor}In [{\color{incolor}3}]:} \PY{l+s+sd}{\PYZdq{}\PYZdq{}\PYZdq{}}
        \PY{l+s+sd}{The initial calculations are done in high\PYZhy{}precision floating point. We convert the}
        \PY{l+s+sd}{results to conventional floating point (float64) later, to speed calculations.}
        \PY{l+s+sd}{\PYZdq{}\PYZdq{}\PYZdq{}}
        \PY{k}{def} \PY{n+nf}{partition\PYZus{}fitness\PYZus{}interval}\PY{p}{(}\PY{n}{min\PYZus{}fitness}\PY{o}{=}\PY{n}{mp\PYZus{}float}\PY{p}{(}\PY{l+s+s1}{\PYZsq{}}\PY{l+s+s1}{\PYZhy{}0.1}\PY{l+s+s1}{\PYZsq{}}\PY{p}{)}\PY{p}{,} 
                                       \PY{n}{max\PYZus{}fitness}\PY{o}{=}\PY{n}{mp\PYZus{}float}\PY{p}{(}\PY{l+s+s1}{\PYZsq{}}\PY{l+s+s1}{0.1}\PY{l+s+s1}{\PYZsq{}}\PY{p}{)}\PY{p}{,}
                                       \PY{n}{subinterval\PYZus{}length}\PY{o}{=}\PY{n}{mp\PYZus{}float}\PY{p}{(}\PY{l+s+s1}{\PYZsq{}}\PY{l+s+s1}{5e\PYZhy{}4}\PY{l+s+s1}{\PYZsq{}}\PY{p}{)}\PY{p}{)}\PY{p}{:}
            \PY{n}{n\PYZus{}fitness\PYZus{}classes} \PY{o}{=} \PY{n+nb}{round}\PY{p}{(}\PY{p}{(}\PY{n}{max\PYZus{}fitness} \PY{o}{\PYZhy{}} \PY{n}{min\PYZus{}fitness}\PY{p}{)} \PY{o}{/} \PY{n}{subinterval\PYZus{}length}\PY{p}{)} \PY{o}{+} \PY{l+m+mi}{1}
            \PY{n}{factors} \PY{o}{=} \PY{n}{Factors}\PY{p}{(}\PY{n}{n\PYZus{}fitness\PYZus{}classes}\PY{p}{,} \PY{n}{death}\PY{o}{=}\PY{o}{\PYZhy{}}\PY{n}{min\PYZus{}fitness}\PY{p}{,} \PY{n}{max\PYZus{}growth}\PY{o}{=}\PY{n}{max\PYZus{}fitness}\PY{p}{)}
            \PY{k}{return} \PY{n}{factors}
        
        \PY{n}{factors} \PY{o}{=} \PY{n}{partition\PYZus{}fitness\PYZus{}interval}\PY{p}{(}\PY{p}{)}
        \PY{n+nb}{print}\PY{p}{(}\PY{l+s+s1}{\PYZsq{}}\PY{l+s+s1}{Number of fitness classes        :}\PY{l+s+s1}{\PYZsq{}}\PY{p}{,} \PY{n}{factors}\PY{o}{.}\PY{n}{n\PYZus{}classes}\PY{p}{)}
        \PY{n+nb}{print}\PY{p}{(}\PY{l+s+s1}{\PYZsq{}}\PY{l+s+s1}{Length of fitness subintervals   :}\PY{l+s+s1}{\PYZsq{}}\PY{p}{,} \PY{n}{factors}\PY{o}{.}\PY{n}{delta}\PY{p}{)}
        \PY{n+nb}{print}\PY{p}{(}\PY{l+s+s1}{\PYZsq{}}\PY{l+s+s1}{Floating\PYZhy{}point precision (digits):}\PY{l+s+s1}{\PYZsq{}}\PY{p}{,} \PY{n}{mp}\PY{o}{.}\PY{n}{dps}\PY{p}{)}
\end{Verbatim}


    \begin{Verbatim}[commandchars=\\\{\}]
Number of fitness classes        : 401
Length of fitness subintervals   : 0.0005
Floating-point precision (digits): 50

    \end{Verbatim}

    \subsubsection{1.2~~Define the discrete distribution of probability over
mutational
effects}\label{define-the-discrete-distribution-of-probability-over-mutational-effects}

Next we discretize a weighted double-gamma distribution. The weight
(mass) of the positive mutational effects in the continuous distribution
is specified by the parameter
\texttt{fraction\_of\_mutations\_that\_increase\_fitness} in the first
cell of this notebook.

    \begin{Verbatim}[commandchars=\\\{\}]
{\color{incolor}In [{\color{incolor}4}]:} \PY{n}{mutations} \PY{o}{=} \PY{n}{WeightedDoubleGamma}\PY{p}{(}\PY{n}{factors}\PY{p}{,} \PY{n}{weight}\PY{o}{=}\PY{n}{mp\PYZus{}float}\PY{p}{(}\PY{n}{fraction\PYZus{}of\PYZus{}mutations\PYZus{}that\PYZus{}increase\PYZus{}fitness}\PY{p}{)}\PY{p}{)}
        \PY{n}{mutations}\PY{o}{.}\PY{n}{normalize}\PY{p}{(}\PY{p}{)}
            
        \PY{n+nb}{print}\PY{p}{(}\PY{l+s+s1}{\PYZsq{}}\PY{l+s+s1}{Parameters of weighted double\PYZhy{}gamma distribution}\PY{l+s+s1}{\PYZsq{}}\PY{p}{)}
        \PY{n+nb}{print}\PY{p}{(}\PY{l+s+s1}{\PYZsq{}}\PY{l+s+s1}{    alpha :}\PY{l+s+s1}{\PYZsq{}}\PY{p}{,} \PY{n}{mutations}\PY{o}{.}\PY{n}{alpha}\PY{p}{)}
        \PY{n+nb}{print}\PY{p}{(}\PY{l+s+s1}{\PYZsq{}}\PY{l+s+s1}{    beta  :}\PY{l+s+s1}{\PYZsq{}}\PY{p}{,} \PY{n}{mutations}\PY{o}{.}\PY{n}{beta}\PY{p}{)}
        \PY{n+nb}{print}\PY{p}{(}\PY{l+s+s1}{\PYZsq{}}\PY{l+s+s1}{    weight:}\PY{l+s+s1}{\PYZsq{}}\PY{p}{,} \PY{n}{mutations}\PY{o}{.}\PY{n}{weight}\PY{p}{)}
        \PY{n+nb}{print}\PY{p}{(}\PY{l+s+s1}{\PYZsq{}}\PY{l+s+s1}{Mean effect of mutation on fitness :}\PY{l+s+s1}{\PYZsq{}}\PY{p}{,} \PY{n}{mutations}\PY{o}{.}\PY{n}{mean}\PY{p}{(}\PY{p}{)}\PY{p}{)}
        \PY{n+nb}{print}\PY{p}{(}\PY{l+s+s1}{\PYZsq{}}\PY{l+s+s1}{Deleterious\PYZhy{}to\PYZhy{}advantageous ratio  :}\PY{l+s+s1}{\PYZsq{}}\PY{p}{,} \PY{n}{mutations}\PY{o}{.}\PY{n}{deleterious\PYZus{}to\PYZus{}advantageous}\PY{p}{(}\PY{p}{)}\PY{p}{)}
        \PY{n+nb}{print}\PY{p}{(}\PY{l+s+s1}{\PYZsq{}}\PY{l+s+s1}{Probability of zero mutation effect:}\PY{l+s+s1}{\PYZsq{}}\PY{p}{,} \PY{n}{mutations}\PY{p}{[}\PY{n}{mutations}\PY{o}{.}\PY{n}{zero\PYZus{}index}\PY{p}{]}\PY{p}{)}
\end{Verbatim}


    \begin{Verbatim}[commandchars=\\\{\}]
Parameters of weighted double-gamma distribution
    alpha : 0.5
    beta  : 500.0
    weight: 0.00000000000001
Mean effect of mutation on fitness : -0.00098322289169388303662041258040983711392852317221381
Deleterious-to-advantageous ratio  : 99999999999999.0
Probability of zero mutation effect: 0.38292492254802620727540922121667547976720435181448

    \end{Verbatim}

    \subsubsection{1.3~~Define the initial frequency distribution of the
population over fitness
classes}\label{define-the-initial-frequency-distribution-of-the-population-over-fitness-classes}

Discretize a zero-mean Gaussian distribution over fitness. The standard
deviation of the continuous distribution is determined by the parameter
\texttt{standard\_deviation\_in\_fitness\_of\_initial\_population},
which is defined in the first cell of this notebook. The mean and
standard deviation of the discrete distribution generally do not match
the parameters of the continuous distribution exactly.

The population is theoretically infinite, but the frequencies of fitness
classes are necessarily finite in the computation. Some solvers of the
initial value problem give more-accurate results when the initial
frequencies are midrange in magnitude, so we scale the initial
frequencies by \(2^{512}.\) The reason for making the scalar an integer
power of 2 is that the operation changes only the exponents, not the
mantissas, in the floating-point representation of the frequencies.

    \begin{Verbatim}[commandchars=\\\{\}]
{\color{incolor}In [{\color{incolor}5}]:} \PY{n}{std} \PY{o}{=} \PY{n}{mp\PYZus{}float}\PY{p}{(}\PY{n}{standard\PYZus{}deviation\PYZus{}in\PYZus{}fitness\PYZus{}of\PYZus{}initial\PYZus{}population}\PY{p}{)}
        
        \PY{l+s+sd}{\PYZdq{}\PYZdq{}\PYZdq{}}
        \PY{l+s+sd}{Here we do not crop the tails of the distribution as BS do. Instead of setting}
        \PY{l+s+sd}{probability masses proportional to probability densities as they do, we integrate}
        \PY{l+s+sd}{the density function over subintervals.}
        \PY{l+s+sd}{\PYZdq{}\PYZdq{}\PYZdq{}}
        \PY{n}{initial\PYZus{}frequencies} \PY{o}{=} \PY{n}{GaussianFrequencies}\PY{p}{(}\PY{n}{factors}\PY{p}{,} \PY{n}{mean}\PY{o}{=}\PY{l+m+mi}{0}\PY{p}{,} \PY{n}{std}\PY{o}{=}\PY{n}{std}\PY{p}{,} \PY{n}{crop}\PY{o}{=}\PY{n}{np}\PY{o}{.}\PY{n}{inf}\PY{p}{,} \PY{n}{density}\PY{o}{=}\PY{k+kc}{False}\PY{p}{)}
        \PY{n}{initial\PYZus{}frequencies}\PY{p}{[}\PY{p}{:}\PY{p}{]} \PY{o}{*}\PY{o}{=} \PY{l+m+mf}{2.0} \PY{o}{*}\PY{o}{*} \PY{l+m+mi}{512}
        \PY{n}{initial\PYZus{}mean}\PY{p}{,} \PY{n}{initial\PYZus{}variance} \PY{o}{=} \PY{n}{initial\PYZus{}frequencies}\PY{o}{.}\PY{n}{mean\PYZus{}and\PYZus{}variance}\PY{p}{(}\PY{p}{)}
        
        
        \PY{n+nb}{print}\PY{p}{(}\PY{l+s+s1}{\PYZsq{}}\PY{l+s+s1}{Fitness of initial population}\PY{l+s+s1}{\PYZsq{}}\PY{p}{)}
        \PY{n+nb}{print}\PY{p}{(}\PY{l+s+s1}{\PYZsq{}}\PY{l+s+s1}{    Mean         :}\PY{l+s+s1}{\PYZsq{}}\PY{p}{,} \PY{n}{initial\PYZus{}mean}\PY{p}{)}
        \PY{n+nb}{print}\PY{p}{(}\PY{l+s+s1}{\PYZsq{}}\PY{l+s+s1}{    Std deviation: }\PY{l+s+s1}{\PYZsq{}}\PY{p}{,} \PY{n}{initial\PYZus{}variance} \PY{o}{*}\PY{o}{*} \PY{l+m+mf}{0.5}\PY{p}{)}
        \PY{n+nb}{print}\PY{p}{(}\PY{l+s+s1}{\PYZsq{}}\PY{l+s+s1}{    Variance     : }\PY{l+s+s1}{\PYZsq{}}\PY{p}{,} \PY{n}{initial\PYZus{}variance}\PY{p}{)}
\end{Verbatim}


    \begin{Verbatim}[commandchars=\\\{\}]
Fitness of initial population
    Mean         : -7.293580011919260841786183682431766807686623890893e-53
    Std deviation:  0.020000380812456165896716463255036527047123607081067
    Variance     :  0.00040001523264326476263976153830639837084670296905497

    \end{Verbatim}

    \subsubsection{1.4~~Define the population, and calculate its equilibrium
distribution}\label{define-the-population-and-calculate-its-equilibrium-distribution}

Our main use of the \texttt{Population} object defined in the following
cell is to calculate derivatives \(\text{d}P_i/\text{d}t\) of the
evolutionary process for a given frequency distribution \(P.\) The
setting \texttt{matrix=True} causes the computation to take the form

\[\mathbf{P^\prime} = W \mathbf{P},\]

as expressed in Section 4 of Basener and Sanford. The equilibrium
distribution of the evolutionary process is a scalar multiple of the
eigenvector of matrix \(W\) corresponding to the largest of its
eigenvalues (again, see Section 4 of Basener and Sanford).

    \begin{Verbatim}[commandchars=\\\{\}]
{\color{incolor}In [{\color{incolor}6}]:} \PY{l+s+sd}{\PYZdq{}\PYZdq{}\PYZdq{}}
        \PY{l+s+sd}{Change from multiprecision to conventional floating point for speed. The statistics}
        \PY{l+s+sd}{reported above will no longer be precisely correct.}
        \PY{l+s+sd}{\PYZdq{}\PYZdq{}\PYZdq{}}
        \PY{n}{factors}\PY{o}{.}\PY{n}{convert}\PY{p}{(}\PY{n+nb}{float}\PY{p}{)}
        \PY{n}{mutations}\PY{o}{.}\PY{n}{convert}\PY{p}{(}\PY{n+nb}{float}\PY{p}{)}  
        \PY{n}{initial\PYZus{}frequencies}\PY{o}{.}\PY{n}{convert}\PY{p}{(}\PY{n+nb}{float}\PY{p}{)}
        
        \PY{n}{pop} \PY{o}{=} \PY{n}{Population}\PY{p}{(}\PY{n}{initial\PYZus{}frequencies}\PY{p}{,} \PY{n}{mutations}\PY{p}{,} \PY{n}{label}\PY{o}{=}\PY{l+s+s1}{\PYZsq{}}\PY{l+s+s1}{ODE solution}\PY{l+s+s1}{\PYZsq{}}\PY{p}{,} \PY{n}{matrix}\PY{o}{=}\PY{k+kc}{True}\PY{p}{)}
        \PY{n}{initial\PYZus{}mean} \PY{o}{=} \PY{n+nb}{float}\PY{p}{(}\PY{n}{pop}\PY{o}{.}\PY{n}{mean}\PY{p}{(}\PY{p}{)}\PY{p}{)}
        \PY{n}{equilibrium} \PY{o}{=} \PY{n}{pop}\PY{o}{.}\PY{n}{equilibrium}\PY{p}{(}\PY{p}{)}
\end{Verbatim}


    \subsection{2~~Verify that variance and correction sum to zero for the
equilibrium
distribution}\label{verify-that-variance-and-correction-sum-to-zero-for-the-equilibrium-distribution}

    As indicated in the introduction, the sum of the variance and the
correction term, which is to say, the rate of change in mean fitness of
the population according to the main theorem of Basener and Sanford,
approaches zero as the evolutionary process converges on the equilibrium
distribution. Of course, we should not expect to see exactly zero in a
numerical calculation of this kind.

    \begin{Verbatim}[commandchars=\\\{\}]
{\color{incolor}In [{\color{incolor}7}]:} \PY{n}{variance}\PY{p}{,} \PY{n}{correction} \PY{o}{=} \PY{n}{pop}\PY{o}{.}\PY{n}{evaluate\PYZus{}theorem}\PY{p}{(}\PY{n}{equilibrium}\PY{p}{)}
        
        \PY{n+nb}{print}\PY{p}{(}\PY{l+s+s1}{\PYZsq{}}\PY{l+s+s1}{Evaluation of Equation 3.4 at the equilibrium distribution}\PY{l+s+s1}{\PYZsq{}}\PY{p}{)}
        \PY{n+nb}{print}\PY{p}{(}\PY{l+s+s1}{\PYZsq{}}\PY{l+s+s1}{    Variance                      : }\PY{l+s+si}{\PYZob{}:+20.17e\PYZcb{}}\PY{l+s+s1}{\PYZsq{}}\PY{o}{.}\PY{n}{format}\PY{p}{(}\PY{n}{variance}\PY{p}{)}\PY{p}{)}
        \PY{n+nb}{print}\PY{p}{(}\PY{l+s+s1}{\PYZsq{}}\PY{l+s+s1}{    Correction                    : }\PY{l+s+si}{\PYZob{}:+20.17e\PYZcb{}}\PY{l+s+s1}{\PYZsq{}}\PY{o}{.}\PY{n}{format}\PY{p}{(}\PY{n}{correction}\PY{p}{)}\PY{p}{)}
        \PY{n+nb}{print}\PY{p}{(}\PY{l+s+s1}{\PYZsq{}}\PY{l+s+s1}{    Rate of change in mean fitness: }\PY{l+s+si}{\PYZob{}:+20.17e\PYZcb{}}\PY{l+s+s1}{\PYZsq{}}\PY{o}{.}\PY{n}{format}\PY{p}{(}\PY{n}{variance}\PY{o}{+}\PY{n}{correction}\PY{p}{)}\PY{p}{)}
\end{Verbatim}


    \begin{Verbatim}[commandchars=\\\{\}]
Evaluation of Equation 3.4 at the equilibrium distribution
    Variance                      : +1.00837812515273403e-04
    Correction                    : -1.00837812515271966e-04
    Rate of change in mean fitness: +1.43656787854329338e-18

    \end{Verbatim}

    \textbf{Establish also that the mean fitness is positive for the
equilibrium distribution.}

    \begin{Verbatim}[commandchars=\\\{\}]
{\color{incolor}In [{\color{incolor}8}]:} \PY{n+nb}{print}\PY{p}{(}\PY{l+s+s1}{\PYZsq{}}\PY{l+s+s1}{Mean fitness over the long term:}\PY{l+s+s1}{\PYZsq{}}\PY{p}{,} \PY{n+nb}{float}\PY{p}{(}\PY{n}{accurate\PYZus{}mean}\PY{p}{(}\PY{n}{equilibrium}\PY{p}{,} \PY{n}{factors}\PY{o}{.}\PY{n}{growth}\PY{p}{)}\PY{p}{)}\PY{p}{)}
\end{Verbatim}


    \begin{Verbatim}[commandchars=\\\{\}]
Mean fitness over the long term: 0.0025584466829794913

    \end{Verbatim}

    \subsection{3~~Check a prediction of the initial change in mean
fitness}\label{check-a-prediction-of-the-initial-change-in-mean-fitness}

Here the sense of \emph{prediction} is that of theoretical prediction of
what we will observe in a numerical solution of the
differential-equation model of the evolutionary process. That is, we use
theory (Equation 3.4) to calculate the instantaneous rate of change in
mean fitness at time 0, and make a linear approximation to the amount of
change that will occur over a short period of time \(x.\) Then we solve
for the frequency distribution at time \(x,\) given the distribution at
time 0, and \emph{observe} (directly calculate) the change in mean
fitness over the time interval \([0, x].\)

Note that this section is primarily pedagogical in purpose. The
calculations performed at time 0 will be performed throughout the
evolutionary process in the next section.

    \subsubsection{3.1~~Calculate the rate of change in mean fitness at time
zero according to Equation
3.4}\label{calculate-the-rate-of-change-in-mean-fitness-at-time-zero-according-to-equation-3.4}

    \begin{Verbatim}[commandchars=\\\{\}]
{\color{incolor}In [{\color{incolor}9}]:} \PY{l+s+sd}{\PYZdq{}\PYZdq{}\PYZdq{}}
        \PY{l+s+sd}{Calculate the initial rate of change according to Eq. 3.4.}
        \PY{l+s+sd}{\PYZdq{}\PYZdq{}\PYZdq{}}
        \PY{n}{variance}\PY{p}{,} \PY{n}{correction} \PY{o}{=} \PY{n}{pop}\PY{o}{.}\PY{n}{evaluate\PYZus{}theorem}\PY{p}{(}\PY{n}{initial\PYZus{}frequencies}\PY{p}{[}\PY{p}{:}\PY{p}{]}\PY{p}{)}
        \PY{n}{rate\PYZus{}of\PYZus{}change\PYZus{}in\PYZus{}mean\PYZus{}fitness} \PY{o}{=} \PY{n}{variance} \PY{o}{+} \PY{n}{correction}
        
        \PY{n+nb}{print}\PY{p}{(}\PY{l+s+s1}{\PYZsq{}}\PY{l+s+s1}{Evaluation of Eq. 3.4 at time 0}\PY{l+s+s1}{\PYZsq{}}\PY{p}{)}
        \PY{n+nb}{print}\PY{p}{(}\PY{l+s+s1}{\PYZsq{}}\PY{l+s+s1}{    Variance                      : }\PY{l+s+si}{\PYZob{}:+20.17e\PYZcb{}}\PY{l+s+s1}{\PYZsq{}}\PY{o}{.}\PY{n}{format}\PY{p}{(}\PY{n}{variance}\PY{p}{)}\PY{p}{)}
        \PY{n+nb}{print}\PY{p}{(}\PY{l+s+s1}{\PYZsq{}}\PY{l+s+s1}{    Correction                    : }\PY{l+s+si}{\PYZob{}:+20.17e\PYZcb{}}\PY{l+s+s1}{\PYZsq{}}\PY{o}{.}\PY{n}{format}\PY{p}{(}\PY{n}{correction}\PY{p}{)}\PY{p}{)}
        \PY{n+nb}{print}\PY{p}{(}\PY{l+s+s1}{\PYZsq{}}\PY{l+s+s1}{    Rate of change in mean fitness: }\PY{l+s+si}{\PYZob{}:+20.17e\PYZcb{}}\PY{l+s+s1}{\PYZsq{}}\PY{o}{.}\PY{n}{format}\PY{p}{(}\PY{n}{rate\PYZus{}of\PYZus{}change\PYZus{}in\PYZus{}mean\PYZus{}fitness}\PY{p}{)}\PY{p}{)}
\end{Verbatim}


    \begin{Verbatim}[commandchars=\\\{\}]
Evaluation of Eq. 3.4 at time 0
    Variance                      : +4.00015232643264771e-04
    Correction                    : -9.83222871491203163e-05
    Rate of change in mean fitness: +3.01692945494144455e-04

    \end{Verbatim}

    \subsubsection{\texorpdfstring{3.2~~Predict the change in mean fitness
over the time interval
\([0,\, x]\)}{3.2~~Predict the change in mean fitness over the time interval {[}0,\textbackslash{}, x{]}}}\label{predict-the-change-in-mean-fitness-over-the-time-interval-0-x}

With rate of change \({\mathrm{d}\bar{m}}/{\mathrm{d}t}\) in mean
fitness at time 0, the mean fitness of the population should change by
approximately \(x \cdot {\mathrm{d}\bar{m}}/{\mathrm{d}t}\) over a short
time interval \([0,\, x].\) We set \(x := 1/128.\)

    \begin{Verbatim}[commandchars=\\\{\}]
{\color{incolor}In [{\color{incolor}10}]:} \PY{n}{x} \PY{o}{=} \PY{l+m+mi}{1}\PY{o}{/}\PY{l+m+mi}{128}
         \PY{l+s+sd}{\PYZdq{}\PYZdq{}\PYZdq{}}
         \PY{l+s+sd}{Multiply the rate of change in mean fitness at time 0, calculated according to }
         \PY{l+s+sd}{Eq. 3.4, by the length x of the time interval [0, x] to approximate the change}
         \PY{l+s+sd}{in mean fitness over the interval.}
         \PY{l+s+sd}{\PYZdq{}\PYZdq{}\PYZdq{}}
         \PY{n}{predicted\PYZus{}change} \PY{o}{=} \PY{n}{x} \PY{o}{*} \PY{n}{rate\PYZus{}of\PYZus{}change\PYZus{}in\PYZus{}mean\PYZus{}fitness}
         \PY{n+nb}{print}\PY{p}{(}\PY{l+s+s1}{\PYZsq{}}\PY{l+s+s1}{Predicted change in mean fitness at time x=1/}\PY{l+s+si}{\PYZob{}0\PYZcb{}}\PY{l+s+s1}{ is }\PY{l+s+si}{\PYZob{}1\PYZcb{}}\PY{l+s+s1}{\PYZsq{}}\PY{o}{.}\PY{n}{format}\PY{p}{(}\PY{n+nb}{int}\PY{p}{(}\PY{l+m+mi}{1}\PY{o}{/}\PY{n}{x}\PY{p}{)}\PY{p}{,} \PY{n}{predicted\PYZus{}change}\PY{p}{)}\PY{p}{)}
\end{Verbatim}


    \begin{Verbatim}[commandchars=\\\{\}]
Predicted change in mean fitness at time x=1/128 is 2.3569761366730036e-06

    \end{Verbatim}

    \subsubsection{3.3~~Compare the prediction to the observed result of
numerical
solution}\label{compare-the-prediction-to-the-observed-result-of-numerical-solution}

    \begin{Verbatim}[commandchars=\\\{\}]
{\color{incolor}In [{\color{incolor}11}]:} \PY{k}{def} \PY{n+nf}{ivp\PYZus{}solution}\PY{p}{(}\PY{n}{derivative}\PY{p}{,} \PY{n}{initial\PYZus{}frequencies}\PY{p}{,} \PY{n}{times}\PY{p}{,} \PY{n}{max\PYZus{}step}\PY{o}{=}\PY{l+m+mi}{1}\PY{o}{/}\PY{l+m+mi}{128}\PY{p}{)}\PY{p}{:}
             \PY{l+s+sd}{\PYZdq{}\PYZdq{}\PYZdq{}}
         \PY{l+s+sd}{    Use the Runge\PYZhy{}Kutta 4(5) method to solve for frequencies at specified `times`, given}
         \PY{l+s+sd}{    the frequencies at the first of those times.}
         \PY{l+s+sd}{    \PYZdq{}\PYZdq{}\PYZdq{}}
             \PY{n}{interval} \PY{o}{=} \PY{p}{(}\PY{n}{times}\PY{p}{[}\PY{l+m+mi}{0}\PY{p}{]}\PY{p}{,} \PY{n}{times}\PY{p}{[}\PY{o}{\PYZhy{}}\PY{l+m+mi}{1}\PY{p}{]} \PY{o}{+} \PY{n}{max\PYZus{}step}\PY{p}{)}
             \PY{n}{result} \PY{o}{=} \PY{n}{solve\PYZus{}ivp}\PY{p}{(}\PY{n}{derivative}\PY{p}{,} \PY{n}{interval}\PY{p}{,} \PY{n}{initial\PYZus{}frequencies}\PY{p}{,} \PY{n}{method}\PY{o}{=}\PY{l+s+s1}{\PYZsq{}}\PY{l+s+s1}{RK45}\PY{l+s+s1}{\PYZsq{}}\PY{p}{,} \PY{n}{t\PYZus{}eval}\PY{o}{=}\PY{n}{times}\PY{p}{,}
                                \PY{n}{rtol}\PY{o}{=}\PY{l+m+mf}{1e\PYZhy{}13}\PY{p}{,} \PY{n}{atol}\PY{o}{=}\PY{l+m+mf}{1e\PYZhy{}11}\PY{p}{,} \PY{n}{max\PYZus{}step}\PY{o}{=}\PY{n}{max\PYZus{}step}\PY{p}{)}
             \PY{k}{return} \PY{n}{result}\PY{o}{.}\PY{n}{y}\PY{o}{.}\PY{n}{T}\PY{p}{,} \PY{n}{result}
         
         \PY{l+s+sd}{\PYZdq{}\PYZdq{}\PYZdq{}}
         \PY{l+s+sd}{Solve for frequencies at time `x`, given frequencies at time 0, and use them to calculate}
         \PY{l+s+sd}{the mean fitness. The solver calls the object `pop` to obtain derivatives.}
         \PY{l+s+sd}{\PYZdq{}\PYZdq{}\PYZdq{}}
         \PY{n}{frequencies}\PY{p}{,} \PY{n}{report} \PY{o}{=} \PY{n}{ivp\PYZus{}solution}\PY{p}{(}\PY{n}{pop}\PY{p}{,} \PY{n}{initial\PYZus{}frequencies}\PY{p}{,} \PY{p}{[}\PY{l+m+mi}{0}\PY{p}{,} \PY{n}{x}\PY{p}{]}\PY{p}{)}
         \PY{n}{observed\PYZus{}change} \PY{o}{=} \PY{n+nb}{float}\PY{p}{(}\PY{n}{accurate\PYZus{}mean}\PY{p}{(}\PY{n}{frequencies}\PY{p}{[}\PY{l+m+mi}{1}\PY{p}{]}\PY{p}{,} \PY{n}{factors}\PY{o}{.}\PY{n}{growth}\PY{p}{)}
                                 \PY{o}{\PYZhy{}} \PY{n}{accurate\PYZus{}mean}\PY{p}{(}\PY{n}{frequencies}\PY{p}{[}\PY{l+m+mi}{0}\PY{p}{]}\PY{p}{,} \PY{n}{factors}\PY{o}{.}\PY{n}{growth}\PY{p}{)}\PY{p}{)}
         
         \PY{n+nb}{print}\PY{p}{(}\PY{l+s+s1}{\PYZsq{}}\PY{l+s+s1}{Status report   :}\PY{l+s+s1}{\PYZsq{}}\PY{p}{,} \PY{n}{report}\PY{o}{.}\PY{n}{message}\PY{p}{)}
         \PY{n+nb}{print}\PY{p}{(}\PY{l+s+s1}{\PYZsq{}}\PY{l+s+s1}{Predicted change:}\PY{l+s+s1}{\PYZsq{}}\PY{p}{,} \PY{n}{predicted\PYZus{}change}\PY{p}{)}
         \PY{n+nb}{print}\PY{p}{(}\PY{l+s+s1}{\PYZsq{}}\PY{l+s+s1}{Observed change :}\PY{l+s+s1}{\PYZsq{}}\PY{p}{,} \PY{n}{observed\PYZus{}change}\PY{p}{)}
         \PY{n+nb}{print}\PY{p}{(}\PY{l+s+s1}{\PYZsq{}}\PY{l+s+s1}{Relative error  :}\PY{l+s+s1}{\PYZsq{}}\PY{p}{,} \PY{p}{(}\PY{n}{predicted\PYZus{}change} \PY{o}{\PYZhy{}} \PY{n}{observed\PYZus{}change}\PY{p}{)} \PY{o}{/} \PY{n}{observed\PYZus{}change}\PY{p}{)}
\end{Verbatim}


    \begin{Verbatim}[commandchars=\\\{\}]
Status report   : The solver successfully reached the end of the integration interval.
Predicted change: 2.3569761366730036e-06
Observed change : 2.356952283365796e-06
Relative error  : 1.0120403105223317e-05

    \end{Verbatim}

    \subsection{4~~Compare predictions to observations throughout an
evolutionary
process}\label{compare-predictions-to-observations-throughout-an-evolutionary-process}

Next we essentially do at 10 thousand times in the evolutionary process
what we just did at the initial time 0. We refer to a number obtained by
evaluation of Equation 3.4 as a \emph{prediction}, and to a number
obtained by direct calculation of the change in mean fitness over a
period of 1/128 year as an \emph{observation.}

    \subsubsection{4.1~~Calculate predictions and
observations}\label{calculate-predictions-and-observations}

    \begin{Verbatim}[commandchars=\\\{\}]
{\color{incolor}In [{\color{incolor}12}]:} \PY{o}{\PYZpc{}\PYZpc{}}\PY{k}{time}
         
         n\PYZus{}years = 10000
         
         \PYZdq{}\PYZdq{}\PYZdq{}
         It is important to set the small time step `dt` to an integer power of 2. This
         ensures that there is no roundoff error when `dt` is added to an integer time,
         provided that the integer is not too large. Similarly, multiplication of the
         calculated rate of change by `dt` changes only the exponent, not the mantissa.
         \PYZdq{}\PYZdq{}\PYZdq{}
         dt = 1/128
         
         \PYZdq{}\PYZdq{}\PYZdq{}
         Set times to 0, 0+dt, 1, 1+dt, ..., n\PYZus{}years, n\PYZus{}years+dt.
         \PYZdq{}\PYZdq{}\PYZdq{}
         times = np.empty(2 * (n\PYZus{}years + 1))
         times[0::2] = np.linspace(0, n\PYZus{}years, n\PYZus{}years+1)  \PYZsh{} set even\PYZhy{}indexed times
         times[1::2] = times[0::2] + dt                  \PYZsh{} set odd\PYZhy{}indexed times
         
         \PYZdq{}\PYZdq{}\PYZdq{}
         Solve for the frequency distribution at each of the times, given the initial
         frequency distribution. The solver calls `pop` to obtain derivatives.
         \PYZdq{}\PYZdq{}\PYZdq{}
         frequencies, report = ivp\PYZus{}solution(pop, initial\PYZus{}frequencies, times, max\PYZus{}step=dt)
         print(\PYZsq{}Report from initial value problem solver\PYZsq{})
         print(\PYZsq{}    Status message        :\PYZsq{}, report.message)
         print(\PYZsq{}    Derivative evaluations:\PYZsq{}, report.nfev)
         
         \PYZdq{}\PYZdq{}\PYZdq{}
         Wrap the sequence of even\PYZhy{}indexed frequency distributions with an `Evolution`
         object, to facilitate animation below.
         \PYZdq{}\PYZdq{}\PYZdq{}
         ev = Evolution(pop)
         ev.set\PYZus{}trajectory(frequencies[::2])
         ev.years\PYZus{}per\PYZus{}epoch = 1
         
         \PYZdq{}\PYZdq{}\PYZdq{}
         Predict the differences in mean fitness at odd\PYZhy{} and even\PYZhy{}indexed times. The
         even\PYZhy{}indexed distributions are used to calculate the rates of change according
         to Eq. 3.4. The calculated rates are multiplied by `dt` to predict the changes
         in mean fitness over intervals of duration `dt`.
         \PYZdq{}\PYZdq{}\PYZdq{}
         theorem\PYZus{}values = np.array([pop.evaluate\PYZus{}theorem(P\PYZus{}t) for P\PYZus{}t in frequencies[0::2]])
         variances = theorem\PYZus{}values[:, 0]
         corrections = theorem\PYZus{}values[:, 1]
         rates\PYZus{}of\PYZus{}change = variances + corrections
         predictions = dt * rates\PYZus{}of\PYZus{}change
         
         \PYZdq{}\PYZdq{}\PYZdq{}
         Calculate the actual differences in mean fitness at odd\PYZhy{} and even\PYZhy{}indexed times.
         The first step is to normalize each of the frequency distributions.
         \PYZdq{}\PYZdq{}\PYZdq{}
         frequencies = (frequencies.T / np.sum(frequencies, axis=1)).T
         means, \PYZus{} = mean\PYZus{}and\PYZus{}variance(factors.growth, frequencies)
         observations = means[1::2] \PYZhy{} means[0::2]
\end{Verbatim}


    \begin{Verbatim}[commandchars=\\\{\}]
Report from initial value problem solver
    Status message        : The solver successfully reached the end of the integration interval.
    Derivative evaluations: 7680020
CPU times: user 13min 44s, sys: 2.42 s, total: 13min 46s
Wall time: 6min 53s

    \end{Verbatim}

    \subsubsection{4.2~~Display results}\label{display-results}

    \begin{Verbatim}[commandchars=\\\{\}]
{\color{incolor}In [{\color{incolor}13}]:} \PY{k+kn}{from} \PY{n+nn}{matplotlib}\PY{n+nn}{.}\PY{n+nn}{ticker} \PY{k}{import} \PY{n}{ScalarFormatter}
         
         \PY{n}{subtitle} \PY{o}{=} \PY{l+s+s1}{\PYZsq{}}\PY{l+s+s1}{Fraction Beneficial }\PY{l+s+si}{\PYZob{}0\PYZcb{}}\PY{l+s+s1}{, Initial Std Dev }\PY{l+s+si}{\PYZob{}1\PYZcb{}}\PY{l+s+s1}{\PYZsq{}}
         \PY{n}{subtitle} \PY{o}{=} \PY{n}{subtitle}\PY{o}{.}\PY{n}{format}\PY{p}{(}\PY{n}{fraction\PYZus{}of\PYZus{}mutations\PYZus{}that\PYZus{}increase\PYZus{}fitness}\PY{p}{,} 
                                    \PY{n}{standard\PYZus{}deviation\PYZus{}in\PYZus{}fitness\PYZus{}of\PYZus{}initial\PYZus{}population}\PY{p}{)}
         
         \PY{k}{def} \PY{n+nf}{animate\PYZus{}evolutionary\PYZus{}process}\PY{p}{(}\PY{n}{ev}\PY{p}{,} \PY{n}{n\PYZus{}years}\PY{p}{,} \PY{n}{n\PYZus{}frames}\PY{o}{=}\PY{l+m+mi}{100}\PY{p}{)}\PY{p}{:}
             \PY{n}{c} \PY{o}{=} \PY{n}{Comparison}\PY{p}{(}\PY{p}{[}\PY{n}{ev}\PY{p}{]}\PY{p}{,} \PY{n}{subtitle}\PY{o}{=}\PY{n}{subtitle}\PY{p}{)}
             \PY{n}{c}\PY{o}{.}\PY{n}{run\PYZus{}until\PYZus{}epoch}\PY{p}{(}\PY{n}{n\PYZus{}years}\PY{p}{)}
             \PY{k}{return} \PY{n}{c}\PY{o}{.}\PY{n}{animate}\PY{p}{(}\PY{n}{nframes}\PY{o}{=}\PY{n}{n\PYZus{}frames}\PY{p}{,} \PY{n}{effective}\PY{o}{=}\PY{k+kc}{False}\PY{p}{)}
             
         \PY{k}{def} \PY{n+nf}{plot\PYZus{}differences}\PY{p}{(}\PY{n}{times}\PY{p}{,} \PY{n}{differences}\PY{p}{)}\PY{p}{:}
             \PY{n}{fig}\PY{p}{,} \PY{n}{ax} \PY{o}{=} \PY{n}{plt}\PY{o}{.}\PY{n}{subplots}\PY{p}{(}\PY{p}{)}
             \PY{n}{ax}\PY{o}{.}\PY{n}{plot}\PY{p}{(}\PY{n}{times}\PY{p}{,} \PY{n}{differences}\PY{p}{)}
             \PY{n}{formatter} \PY{o}{=} \PY{n}{ScalarFormatter}\PY{p}{(}\PY{n}{useMathText}\PY{o}{=}\PY{k+kc}{True}\PY{p}{)}
             \PY{n}{formatter}\PY{o}{.}\PY{n}{set\PYZus{}powerlimits}\PY{p}{(}\PY{p}{(}\PY{o}{\PYZhy{}}\PY{l+m+mi}{1}\PY{p}{,} \PY{l+m+mi}{2}\PY{p}{)}\PY{p}{)}
             \PY{n}{ax}\PY{o}{.}\PY{n}{yaxis}\PY{o}{.}\PY{n}{set\PYZus{}major\PYZus{}formatter}\PY{p}{(}\PY{n}{formatter}\PY{p}{)}
             \PY{n}{ax}\PY{o}{.}\PY{n}{set\PYZus{}ylabel}\PY{p}{(}\PY{l+s+s1}{\PYZsq{}}\PY{l+s+s1}{Error = Predicted – Observed}\PY{l+s+s1}{\PYZsq{}}\PY{p}{)}
             \PY{n}{ax}\PY{o}{.}\PY{n}{set\PYZus{}xlabel}\PY{p}{(}\PY{l+s+s1}{\PYZsq{}}\PY{l+s+s1}{Year}\PY{l+s+s1}{\PYZsq{}}\PY{p}{)}
             \PY{n}{fig}\PY{o}{.}\PY{n}{suptitle}\PY{p}{(}\PY{l+s+s1}{\PYZsq{}}\PY{l+s+s1}{Errors in Predicting Mean Fitness 1/}\PY{l+s+si}{\PYZob{}\PYZcb{}}\PY{l+s+s1}{ Year in Future}\PY{l+s+s1}{\PYZsq{}}\PY{o}{.}\PY{n}{format}\PY{p}{(}\PY{n+nb}{round}\PY{p}{(}\PY{l+m+mi}{1}\PY{o}{/}\PY{n}{dt}\PY{p}{)}\PY{p}{)}\PY{p}{)}
             \PY{n}{ax}\PY{o}{.}\PY{n}{set\PYZus{}title}\PY{p}{(}\PY{n}{subtitle}\PY{p}{)}
             \PY{k}{return} \PY{n}{fig}\PY{p}{,} \PY{n}{ax}
         
         \PY{k}{def} \PY{n+nf}{plot\PYZus{}theorem\PYZus{}values}\PY{p}{(}\PY{n}{times}\PY{p}{,} \PY{n}{variances}\PY{p}{,} \PY{n}{corrections}\PY{p}{,} \PY{n}{observed}\PY{p}{)}\PY{p}{:}
             \PY{n}{fig}\PY{p}{,} \PY{n}{ax} \PY{o}{=} \PY{n}{plt}\PY{o}{.}\PY{n}{subplots}\PY{p}{(}\PY{p}{)}
             \PY{n}{ax}\PY{o}{.}\PY{n}{plot}\PY{p}{(}\PY{n}{times}\PY{p}{,} \PY{n}{variances}\PY{p}{,} \PY{n}{label}\PY{o}{=}\PY{l+s+s1}{\PYZsq{}}\PY{l+s+s1}{Variance}\PY{l+s+s1}{\PYZsq{}}\PY{p}{)}
             \PY{n}{ax}\PY{o}{.}\PY{n}{plot}\PY{p}{(}\PY{n}{times}\PY{p}{,} \PY{n}{corrections}\PY{p}{,} \PY{n}{label}\PY{o}{=}\PY{l+s+s1}{\PYZsq{}}\PY{l+s+s1}{Correction}\PY{l+s+s1}{\PYZsq{}}\PY{p}{)}
             \PY{n}{ax}\PY{o}{.}\PY{n}{plot}\PY{p}{(}\PY{n}{times}\PY{p}{,} \PY{n}{variances} \PY{o}{+} \PY{n}{corrections}\PY{p}{,} \PY{n}{label}\PY{o}{=}\PY{l+s+s1}{\PYZsq{}}\PY{l+s+s1}{Predicted}\PY{l+s+s1}{\PYZsq{}}\PY{p}{)}
             \PY{n}{ax}\PY{o}{.}\PY{n}{plot}\PY{p}{(}\PY{n}{times}\PY{p}{,} \PY{n}{observed}\PY{p}{,} \PY{n}{label}\PY{o}{=}\PY{l+s+s1}{\PYZsq{}}\PY{l+s+s1}{Observed}\PY{l+s+s1}{\PYZsq{}}\PY{p}{,} \PY{n}{ls}\PY{o}{=}\PY{l+s+s1}{\PYZsq{}}\PY{l+s+s1}{\PYZhy{}\PYZhy{}}\PY{l+s+s1}{\PYZsq{}}\PY{p}{)}
             \PY{n}{formatter} \PY{o}{=} \PY{n}{ScalarFormatter}\PY{p}{(}\PY{n}{useMathText}\PY{o}{=}\PY{k+kc}{True}\PY{p}{)}
             \PY{n}{formatter}\PY{o}{.}\PY{n}{set\PYZus{}powerlimits}\PY{p}{(}\PY{p}{(}\PY{o}{\PYZhy{}}\PY{l+m+mi}{1}\PY{p}{,} \PY{l+m+mi}{2}\PY{p}{)}\PY{p}{)}
             \PY{n}{ax}\PY{o}{.}\PY{n}{yaxis}\PY{o}{.}\PY{n}{set\PYZus{}major\PYZus{}formatter}\PY{p}{(}\PY{n}{formatter}\PY{p}{)}
             \PY{n}{ax}\PY{o}{.}\PY{n}{set\PYZus{}xlabel}\PY{p}{(}\PY{l+s+s1}{\PYZsq{}}\PY{l+s+s1}{Year}\PY{l+s+s1}{\PYZsq{}}\PY{p}{)}
             \PY{n}{ax}\PY{o}{.}\PY{n}{set\PYZus{}ylabel}\PY{p}{(}\PY{l+s+s1}{\PYZsq{}}\PY{l+s+s1}{Contribution to Rate of Change in Mean Fitness}\PY{l+s+s1}{\PYZsq{}}\PY{p}{)}
             \PY{n}{ax}\PY{o}{.}\PY{n}{legend}\PY{p}{(}\PY{n}{loc}\PY{o}{=}\PY{l+s+s1}{\PYZsq{}}\PY{l+s+s1}{best}\PY{l+s+s1}{\PYZsq{}}\PY{p}{,} \PY{n}{ncol}\PY{o}{=}\PY{l+m+mi}{2}\PY{p}{)}
             \PY{n}{fig}\PY{o}{.}\PY{n}{suptitle}\PY{p}{(}\PY{l+s+s1}{\PYZsq{}}\PY{l+s+s1}{Predicted and Observed Rates of Change in Mean Fitness}\PY{l+s+s1}{\PYZsq{}}\PY{p}{)}
             \PY{n}{ax}\PY{o}{.}\PY{n}{set\PYZus{}title}\PY{p}{(}\PY{n}{subtitle}\PY{p}{)}
             \PY{k}{return} \PY{n}{fig}\PY{p}{,} \PY{n}{ax}
         
         \PY{n}{animate\PYZus{}evolutionary\PYZus{}process}\PY{p}{(}\PY{n}{ev}\PY{p}{,} \PY{n}{n\PYZus{}years}\PY{p}{)}
\end{Verbatim}


\begin{Verbatim}[commandchars=\\\{\}]
{\color{outcolor}Out[{\color{outcolor}13}]:} <matplotlib.animation.FuncAnimation at 0x7f45000c0a20>
\end{Verbatim}
            
    There are "damped oscillations" (for want of a better term) about the
equilibrium distribution as the process converges. The closer the mean
fitness at equilibrium is to zero, the more pronounced are the
oscillations. The wavy lines in the following figure correspond to the
oscillations in the aniation.

    \begin{Verbatim}[commandchars=\\\{\}]
{\color{incolor}In [{\color{incolor}14}]:} \PY{n}{plot\PYZus{}theorem\PYZus{}values}\PY{p}{(}\PY{n}{times}\PY{p}{[}\PY{l+m+mi}{0}\PY{p}{:}\PY{p}{:}\PY{l+m+mi}{2}\PY{p}{]}\PY{p}{,} \PY{n}{variances}\PY{p}{,} \PY{n}{corrections}\PY{p}{,} \PY{n}{observations} \PY{o}{/} \PY{n}{dt}\PY{p}{)}
         \PY{n}{plot\PYZus{}differences}\PY{p}{(}\PY{n}{times}\PY{p}{[}\PY{l+m+mi}{0}\PY{p}{:}\PY{p}{:}\PY{l+m+mi}{2}\PY{p}{]}\PY{p}{,} \PY{p}{(}\PY{n}{predictions} \PY{o}{\PYZhy{}} \PY{n}{observations}\PY{p}{)}\PY{p}{)}\PY{p}{;}
\end{Verbatim}


    
    \begin{verbatim}
<IPython.core.display.Javascript object>
    \end{verbatim}

    
    
    \begin{verbatim}
<IPython.core.display.HTML object>
    \end{verbatim}

    
    
    \begin{verbatim}
<IPython.core.display.Javascript object>
    \end{verbatim}

    
    
    \begin{verbatim}
<IPython.core.display.HTML object>
    \end{verbatim}

    
    \subsection{5~~Compare the final frequency distribution to the
equilibrium
distribution}\label{compare-the-final-frequency-distribution-to-the-equilibrium-distribution}

    \begin{Verbatim}[commandchars=\\\{\}]
{\color{incolor}In [{\color{incolor}15}]:} \PY{k}{def} \PY{n+nf}{plot\PYZus{}equilibrium}\PY{p}{(}\PY{n}{equilibrium}\PY{p}{,} \PY{n}{final\PYZus{}frequencies}\PY{p}{)}\PY{p}{:}
             \PY{n}{fig}\PY{p}{,} \PY{n}{ax} \PY{o}{=} \PY{n}{plt}\PY{o}{.}\PY{n}{subplots}\PY{p}{(}\PY{p}{)}
             \PY{n}{fig}\PY{o}{.}\PY{n}{suptitle}\PY{p}{(}\PY{l+s+s1}{\PYZsq{}}\PY{l+s+s1}{Comparison of Equilibrium Distribution and Final Frequencies}\PY{l+s+s1}{\PYZsq{}}\PY{p}{)}
             \PY{n}{ax}\PY{o}{.}\PY{n}{set\PYZus{}title}\PY{p}{(}\PY{n}{subtitle}\PY{p}{)}
             \PY{n}{ax}\PY{o}{.}\PY{n}{set\PYZus{}xlabel}\PY{p}{(}\PY{l+s+s1}{\PYZsq{}}\PY{l+s+s1}{Fitness}\PY{l+s+s1}{\PYZsq{}}\PY{p}{)}
             \PY{n}{ax}\PY{o}{.}\PY{n}{set\PYZus{}ylabel}\PY{p}{(}\PY{l+s+s1}{\PYZsq{}}\PY{l+s+s1}{Relative Frequency}\PY{l+s+s1}{\PYZsq{}}\PY{p}{)}
             \PY{n}{plt}\PY{o}{.}\PY{n}{plot}\PY{p}{(}\PY{n}{factors}\PY{o}{.}\PY{n}{growth}\PY{p}{,} \PY{n}{equilibrium}\PY{p}{,} \PY{n}{label}\PY{o}{=}\PY{l+s+s1}{\PYZsq{}}\PY{l+s+s1}{Equilibrium}\PY{l+s+s1}{\PYZsq{}}\PY{p}{)}
             \PY{n}{plt}\PY{o}{.}\PY{n}{plot}\PY{p}{(}\PY{n}{factors}\PY{o}{.}\PY{n}{growth}\PY{p}{,} \PY{n}{final\PYZus{}frequencies}\PY{p}{,} \PY{n}{label}\PY{o}{=}\PY{l+s+s1}{\PYZsq{}}\PY{l+s+s1}{Last Distribution}\PY{l+s+s1}{\PYZsq{}}\PY{p}{,} \PY{n}{ls}\PY{o}{=}\PY{l+s+s1}{\PYZsq{}}\PY{l+s+s1}{\PYZhy{}\PYZhy{}}\PY{l+s+s1}{\PYZsq{}}\PY{p}{)}
             \PY{n}{plt}\PY{o}{.}\PY{n}{legend}\PY{p}{(}\PY{n}{loc}\PY{o}{=}\PY{l+s+s1}{\PYZsq{}}\PY{l+s+s1}{best}\PY{l+s+s1}{\PYZsq{}}\PY{p}{)}
             \PY{k}{return} \PY{n}{fig}\PY{p}{,} \PY{n}{ax}
         
         \PY{n}{plot\PYZus{}equilibrium}\PY{p}{(}\PY{n}{equilibrium}\PY{p}{,} \PY{n}{frequencies}\PY{p}{[}\PY{o}{\PYZhy{}}\PY{l+m+mi}{1}\PY{p}{]}\PY{p}{)}\PY{p}{;}
\end{Verbatim}


    
    \begin{verbatim}
<IPython.core.display.Javascript object>
    \end{verbatim}

    
    
    \begin{verbatim}
<IPython.core.display.HTML object>
    \end{verbatim}

    
    In some cases, 10 thousand years is not sufficient time for convergence,
and the frequency distribution at the end of the process is not a
precise match for the equilibrium distribution in the figure above.


    % Add a bibliography block to the postdoc
    
    
    
    \end{document}
